% 第四章文字初稿
% 编译方式：pdfLaTeX + ctex
% 风格：许居衍院士
\documentclass[11pt,a4paper]{article}
\usepackage[UTF8]{ctex}
\usepackage[top=2cm, bottom=2cm, left=2cm, right=2cm]{geometry}
\usepackage{booktabs}
\usepackage{array}
\usepackage{tabularx}
\usepackage{enumitem}
\usepackage[font=small,skip=4pt]{caption}

\usepackage{indentfirst}
\setlength{\parskip}{2pt}
\setlength{\parindent}{2em}
\setlist{nosep,leftmargin=*}
\renewcommand{\arraystretch}{1.15}

\title{\bfseries\large 第四章\quad 产业技术发展趋势展望}
\author{}
\date{}

\begin{document}

\maketitle

\section{领域专用化成为主流}

半导体产业在过去半个多世纪中，始终沿着一条"统一芯片"（unified chip）的路径演进。这一路径的核心逻辑是：用单一的CMOS工艺平台、冯·诺依曼架构和摩尔定律的节奏，去满足尽可能多的应用需求。CPU是如此，存储器是如此，乃至SoC亦如此。然而，当摩尔定律在2014—2017年间逼近S曲线拐点、登纳德等效缩放（Dennard scaling）早已失效之时，这一"统一"逻辑开始松动。一个根本性的问题浮出水面：当通用平台无法再以可承受的成本提供性能增益时，产业将走向何方？

答案是——领域专用化（domain-specific architecture, DSA）。

所谓领域专用化，是指芯片架构不再追求"万能"，而是针对特定应用领域的计算特征进行深度定制。这不是一个全新的概念，ASIC早在1980年代就已存在。但过去十年间发生的变化是质的：领域专用化从产业边缘的补位角色，跃升为技术创新的主流路径。历史事件犹如枝上嫩芽，总在它要长出的地方露头——2015年Google TPU的发布，就是这枚嫩芽。

2015年，Google公开了其部署在数据中心中的张量处理单元（Tensor Processing Unit，TPU）。这是一款为神经网络推理深度定制的ASIC，基于256×256脉动阵列（systolic array）架构，在75W功耗下提供92 TOPS的推理性能。与同时期CPU和GPU相比，TPU在能效上实现了数量级的提升。该事件标志着云侧AI加速时代的开启。

此后，领域专用加速器在云端呈现百花齐放之势。NVIDIA在其GPU中引入专用张量核心（Tensor Core），将4×4×4矩阵乘累加指令直接硬化，使Volta V100的混合精度训练性能达到120 Tensor TFLOPS，这一思路的本质是"在通用平台上嵌入专用逻辑"。Amazon推出定制AI训练芯片Trainium，单芯片算力达1.3 PFLOPS，并通过NeuronLink实现64芯片互联。Graphcore的智能处理单元（Intelligence Processing Unit，IPU）采用大规模并行BSP执行模型，以1216个独立核心和300MB片上SRAM重构了AI计算的硬件架构。Groq的语言处理单元（Language Processing Unit，LPU）则彻底抛弃缓存一致性模型，以编译器调度的"张量流水线"实现确定性执行——230MB全片上SRAM、80TB/s片上带宽，无片外HBM。值得注意的是，这些架构设计各不相同、甚至截然相反，但有一点是共同的：它们都放弃了CPU式的"万能"追求，转而寻求特定计算模式下的极致效率。这正是领域专用化的本质特征。

「注1」：Jouppi et al., "In-Datacenter Performance Analysis of a Tensor Processing Unit," ISCA 2017。该文是领域专用架构领域引用最高的文献之一。

在移动端，Apple于2017年在A11 Bionic芯片中首次集成专用双核神经网络引擎（Neural Engine），实现6000亿次/秒的推理性能，使Face ID和AR实时处理得以在手机端完成。此后从M1到M5系列，Apple持续强化这一方向：不仅集成了CPU、GPU和Neural Engine，更通过统一内存池（Unified Memory Architecture，UMA）将各种专用加速器之间的数据搬运延迟降至几乎为零，带宽推至800+ GB/s。这一架构的本质是"异构但不分割"——让每个专用模块做自己最擅长的事，同时通过统一的存储和调度机制消除模块间的摩擦。

Mobileye EyeQ5自动驾驶SoC则是领域专用化在汽车电子中的深入体现。这款7nm芯片集成向量微码处理器与可编程宏阵列，在24 TOPS算力下支撑L4/L5级自动驾驶的感知与决策。与通用GPU方案相比，EyeQ5在功耗、实时性和功能安全方面具有不可替代的优势——当生命安全成为首要约束时，"专用"不再是效率选择，而是必然要求。

Cerebras的晶圆级引擎（Wafer-Scale Engine，WSE）代表了领域专用化的另一种极端路径。WSE-3将整片晶圆（46,225mm²）作为单一处理器，不切割、不封装，利用原本废弃的划片线（scribe line）作为die间互连通道。900,000个处理核心通过214 Pb/s的片上网络连接，将深度学习训练的通信瓶颈降为零。这一设计颠覆了半导体制造"先切割再封装"的基本假设——当通用范式无法满足需求时，产业不惜彻底重塑芯片的物理形态。

「注2」：Cerebras的晶圆级方案与chiplet路径形成鲜明对比。前者走向"极端集成"，后者走向"极端分割"，但两者指向同一结论：通用架构已不足以支撑AI时代的计算需求。

领域专用化不仅是技术路线的调整，更是半导体产业价值链的深刻重塑。在通用芯片时代，IDM（集成器件制造）和Fabless（无晶圆设计）是产业的主角，设计、制造、封装各环节相对独立。芯片设计公司只需遵循既定的工艺和架构规则，专注于摩尔定律驱动的微缩路线即可。

但领域专用化改变了这一格局。当芯片架构需要为特定应用领域深度定制时，只有深刻理解应用场景需求的企业才能定义最优的硬件架构。于是，拥有核心应用场景的系统企业开始亲自下场设计芯片。Google的TPU、Amazon的Trainium、Apple的Neural Engine、Tesla的Dojo——这些案例的共同特征是，它们不再满足于购买货架上的通用芯片，而是从系统需求出发，反向定义工艺、架构和封装方案。这一转变意味着"芯片设计"已从专业分工中的一个可外包环节，上升为系统级创新的核心战略能力。数年之前，笔者注意到苹果公司开始收购半导体公司时就已觉察到这一趋势：未来将是大型系统公司介入芯片的时代，如今这一判断已全面应验。

与此同时，AMD Xilinx Versal ACAP代表了另一种思路：将ARM标量引擎、FPGA可编程逻辑和VLIW SIMD AI引擎阵列集成在统一平台上，通过可重构性来弥合通用与专用之间的鸿沟。这种"用硬件可编程性来应对应用多样性"的路径，是领域专用化浪潮中的一个重要分支。

如果将视角拉回到本文的核心论点——半导体循环律的终结——领域专用化的意义便更加清晰。

传统半导体产业的发展，建立在"统一芯片"范式之上：一种工艺（CMOS）、一种架构（冯·诺依曼）、一条演进路径（摩尔定律）。这一范式之所以能够长期有效，是因为在那个阶段，将更多晶体管集成到一颗芯片上就意味着更高的性能和更低的成本。产业运行的逻辑是线性的、可预测的。

然而，当微缩的经济性递减、架构多样化发散、应用需求分岔之后，"统一"的红利被消耗殆尽。这就像一条S曲线走过了快速增长期，进入了平台期——无论再增加多少投入，回报都在递减。而领域专用化，就是在这个转折点上涌现出的多条新S曲线，每一条都面向特定的应用空间，互不重叠、也不可替代。正如笔者在《集成电路发展哲理》中所言："当它在缩小进程中变得愈来愈繁而又不能化简时，就意味着基于CMOS结构在逻辑上不能延伸了。"

「注3」：AMD Chiplet架构（Zen系列）将计算核心和I/O分离至不同制程节点，通过Infinity Fabric互联——其本质是：不再试图用一颗die满足所有需求，而是让每一部分都用最适合自己的方式去实现。这是"统一"逻辑向"多元"逻辑转变的微观缩影。

领域专用化的兴起，是半导体循环律终结在技术和产业层面的最直接体现。从Google TPU到Apple Neural Engine，从Cerebras WSE到Tesla Dojo，这些案例共同标示着同一个方向：通用芯片的时代正在让位于领域专用的时代。

但需要指出的是，领域专用化并非对通用化的简单否定。如果从笔者提出的"由简入繁、化繁为简"的螺旋式发展史观来看，这不过是产业演进中的一个必经阶段。当通用架构在微缩驱动下经历了数十年的"由简入繁"式发展后，如今正进入"化繁为简"的再抽象过程，而领域专用化就是这个过程中的第一次分化。未来，当各条领域专用路线各自发展到一定阶段后，产业可能再度走向新的通用化。这不简单是技术问题，而是半导体产业发展内在逻辑的必然。

\section{异构计算成为核心架构}

如果说领域专用化回答了"为谁而造"的问题，那么异构计算（heterogeneous computing）则回答了"如何组织"的问题——在单一芯片或单一系统内，将不同计算范式、不同架构特点的处理单元组合起来，使每种计算任务都由最适合它的硬件来执行。

异构计算并非新鲜事物。早在2005年，IBM Cell处理器就将一个PowerPC控制核与八个数据并行协处理器整合在一起，开创了异构多核的先河。但真正使其从"特例"变为"主流"的，是后摩尔时代统一芯片范式失效后产业被迫采取的架构多元化策略。当CMOS微缩不再能自动提升性能时，唯一的选择就是让专用硬件做专用的事——而这正是异构计算的核心逻辑。

异构计算最广泛的应用形态，是在单一SoC中集成CPU、GPU和NPU。Apple M系列芯片是这一方向的标杆：它将高性能CPU核心、高能效CPU核心、大规模并行GPU和专用Neural Engine集成在单一芯片上，通过统一内存池（UMA）实现所有处理器对同一数据的零拷贝访问。M3 Ultra的UMA带宽达到800+ GB/s，使得CPU、GPU和NPU之间的数据搬运几乎不存在瓶颈。这一设计的本质，是"让每一个计算单元都做自己最擅长的事"——顺序计算交给CPU，并行渲染交给GPU，神经网络推理交给Neural Engine。

AMD的MI300A APU则将这一思路推向了新的高度。它将24核Zen4 CPU、228组CDNA3 GPU和128GB统一HBM3内存集成在单一封装内，通过Infinity Fabric实现缓存一致性互联。搭载44,564颗MI300A的El Capitan超算登顶Top500榜首，证明了异构融合架构在HPC领域的统治力。

「注4」：AMD的Infinity Fabric是一种芯片间缓存一致性互联技术，能够实现CPU和GPU之间的低延迟数据共享。其核心特点是从硬件层面解决了异构计算中最棘手的"数据同步"问题。

在CPU领域，ARM于2011年率先提出big.LITTLE架构，将高性能大核与高能效小核组合在同一芯片中。这一设计在移动端取得了巨大成功，其后于2017年升级为DynamIQ，允许灵活配置1+3、1+3+4等多种核心组合。Intel于2021年在Alder Lake中引入了类似的混合x86架构，通过Performance-core（P核）和Efficiency-core（E核）的组合，辅以Thread Director硬件调度，实现了性能和功耗的动态平衡。这一趋势表明，即使在同一指令集架构内部，异构化也已从"可选"变为"必需"。

异构计算的另一重要分支是神经形态计算（neuromorphic computing）。与CPU+GPU的异构不同，神经形态芯片从根本上背离了冯·诺依曼架构——它不区分计算和存储，而是通过模拟生物神经元的脉冲发放机制来实现信息处理。

IBM TrueNorth（2014）是这一方向的先驱。这款芯片集成了4096个神经突触核心，包含100万个可编程神经元和2.56亿个可编程突触，在28nm工艺下功耗仅为70mW。其核心设计思想是事件驱动（event-driven）：只有当脉冲到来时，相应的电路才被激活，从而实现了极致的能效。

Intel Loihi 2则将神经形态计算推向了新高度。128个全异步神经形态核心每个都包含192KB SRAM，支持微码可编程的神经元模型和片上学习规则。2024年，Intel发布了Hala Point系统，将1152颗Loihi 2芯片连接在一起，构建了拥有11.5亿神经元的全球最大神经形态系统，功耗仅为2600W。

IBM NorthPole（2023）则代表了存算融合的另一种路径。它将计算和SRAM交织排布，256个核心共192MB片上SRAM，彻底消除了片外DRAM访问。在ResNet-50推理任务上，NorthPole的能效达到同类GPU的25倍。

「注5」：Modha et al., "A 12nm 256-core brain-inspired chip for efficient AI inference," Science, 2023。该文指出NorthPole通过完全消除片外内存访问，实现了计算能效的范式级突破。

存内计算（compute-in-memory）是异构计算中另一条重要的技术路径。它的核心思想是：不再在存储器和处理器之间搬运数据，而是直接在存储阵列中完成计算。

IBM于2023年发布的混合信号存内计算芯片，采用相变存储器（PCM）交叉阵列在模拟域内完成矩阵乘法，再通过片上的ADC/DAC和数字电路完成激活函数的计算。64个核心在14nm CMOS工艺上实现了超过15倍于此前同类芯片的计算密度。

Samsung在2022年发表的MRAM存内计算芯片则改写了游戏规则。传统的存内计算基于RRAM或PCM，但MRAM有着更高的速度和更好的耐久性。Samsung的研究团队发明了一种"电阻-求和"架构，克服了MRAM单元低电阻带来的技术困难，在28nm工艺上实现了98\%的MNIST识别精度。

「注6」：Jung et al., "A crossbar array of magnetoresistive memory devices for in-memory computing," Nature, 2022。该文首次证明了MRAM可用于存内计算，为高耐久、高速度的存内计算打开了新的材料路线。

异构计算的兴起，对半导体产业的制造和设计模式产生了深远影响。在通用芯片时代，一颗芯片的设计和制造可以完全由一家IDM或Fabless+Foundry的组合完成。但在异构计算时代，不同计算单元往往需要不同的工艺节点——CPU追求高性能逻辑工艺，DRAM追求高密度存储工艺，模拟器件需要特殊工艺——将它们集成在一起，就需要先进的封装技术和复杂的chiplet互连方案。

这一趋势催生了新的产业协作模式。AMD的chiplet架构让CPU核心和I/O单元使用不同的制程节点，通过Infinity Fabric互连；NVIDIA的Grace Hopper则将Arm CPU和GPU通过NVLink-C2C实现缓存一致性互联。这些实践表明，异构计算正在推动半导体产业从"单一芯片设计"走向"多芯片系统设计"，从"工艺集成"走向"系统集成"。

异构计算的本质，是对半导体循环律中"统一芯片"假设的根本性突破。在统一范式时代，产业深信"一种架构可以满足所有需求"；而在异构计算时代，产业开始接受一个基本事实：没有一种计算范式可以同时在通用性、性能、能效和成本上达到最优。

由此带来的结果是：产业从"寻找最好的架构"转向"让多种架构协同工作"。这正是笔者关于后循环时代论述的核心判断之一。我们在2008年就指出，"用与造的对立统一律是硅产业发展的一个主线索"，而异构计算正是这一对立统一律在架构层面的具体体现——当通用不再够用时，异构便成为必然选择。

异构计算成为核心架构，是半导体循环律终结后技术范式转型的重要标志。从CPU+GPU的简单融合，到神经形态、存内计算等非冯架构的涌现，异构计算的内涵在不断扩展。但硅-冯范式仍将长期主导产业发展——异构计算的真正意义，并非用新架构取代旧架构，而是让多种架构在统一系统中各司其职、协同工作。这是后循环时代"系统协同"逻辑在架构层面的第一次展开。

\section{混合计算成为重要趋势}

当领域专用化和异构计算在数字芯片范畴内不断深化时，另一条更加根本的技术路径也在加速演进——这就是混合计算（hybrid/novel computing）。与异构计算在数字框架内进行架构创新不同，混合计算试图突破"基于电子电荷的二进制数字计算"这一硅-冯范式的底层假设，探索基于其他物理机制的信息处理方式。

笔者在2021年的演讲中曾指出，器件包括很多种——有电子的、光子的、离子的、激子的——"实际上，所有的物理粒子都可以用来做器件。"这暗示了一个重要的方向：后摩尔时代的计算创新，不仅发生在架构层面，也发生在物理基础层面。

量子计算（quantum computing）利用量子叠加和纠缠等物理特性，在特定计算问题上展现出超越经典计算机的潜力。近年来，这一领域正在从实验室走向工程实现。

Google于2024年发布的Willow量子处理器是这一进展的里程碑。这款105量子比特的超导量子处理器首次实现了"低于阈值"的量子纠错——这是量子计算领域一个长期追求的目标。通过表面码（surface code）技术，Willow在距离3/5/7的纠错实验中展示了指数级的误差抑制，标志着量子纠错从理论走向实践。

在离子阱路线上，Quantinuum的H2处理器达到了56量子比特和2^21的量子体积——这是迄今公开报道的最高值。其全联通拓扑意味着任意两个量子比特之间都可以直接执行两量子门操作，这在量子算法实现中具有重要优势。

Intel的Tunnel Falls芯片则代表了另一条独特的路径——硅自旋量子比特。与超导或离子阱方案不同，Intel采用标准CMOS工艺在300mm晶圆上制造量子点器件，良率达到95\%。这种"借用现有半导体产线"的思路，为量子计算的规模化制造提供了极具吸引力的前景。

「注7」：Intel Tunnel Falls芯片的制造采用了与标准CMOS兼容的工艺，这意味著量子比特的制造可以实现与经典芯片同样的规模经济效应。

光计算（optical/photonic computing）试图用光子替代电子来执行计算，利用光的速度、带宽和低功耗特性突破电子计算的物理限制。

2022年，加拿大的Xanadu公司发布了Borealis光量子处理器。这款基于216个压缩态（squeezed state）的可编程光子处理器，在36微秒内完成了一项估计需要9000年才能在经典超算上完成的任务。Borealis是世界上首个具有完全可编程动态门的光量子系统，验证了光量子计算的技术可行性。

在经典光计算方向，Lightmatter于2025年发布了Passage M1000光子互连芯片，采用3D光子中介层架构，实现了114 Tbps的片间光通信带宽。这一技术有望解决AI训练中日益严重的互连瓶颈——当芯片内计算速度远超片间通信速度时，光子互连提供了一个超越电子互连的解决方案。

清华大学发布的Taichi/Taichi-II光计算芯片则代表了另一条路径：直接在光域内完成神经网络的推理和训练。Taichi芯片实现了160 TOPS/W的能效，而Taichi-II更是首次实现了"全前向模式训练"——通过利用光子传播的对称性，彻底摆脱了训练过程中对GPU的依赖。这一突破使光神经网络的训练速度提升了10倍。

「注8」：Taichi-II的训练方法利用光的物理特性——光子传播的互易性（reciprocity）——来替代传统反向传播算法中的梯度计算，这是"用物理定律替代算法"的典型案例。

除了纯量子计算和光计算之外，还有一条相对务实的中间路径——量子启发式计算（quantum-inspired computing）。这类系统不依赖量子物理效应，但借鉴了量子算法的思想，在经典硬件上实现针对特定优化问题的大幅加速。

Toshiba的模拟分岔机（Simulated Bifurcation Machine，SBM）是这一方向的典型代表。这款基于FPGA的量子启发式Ising机，拥有4096个全连接自旋，能效达到相干Ising机的288倍。在5G基站资源分配等实际应用中，SBM能够在0.5毫秒内完成传统算法需要数秒才能解决的优化问题。这类"无需量子硬件，却能获得量子级加速"的方案，在当前量子计算基础设施尚不完善的阶段，具有重要的实用价值。

混合计算的技术路径虽然多样，但它们的产业化进程却展现出一些共性特征：以云计算为入口、以专用加速为切入、以API（应用程序编程接口）为纽带。目前，主要的量子计算服务商（IBM、Google、Amazon）均已通过云平台提供量子计算接入，用户无需自建量子硬件即可进行算法开发和验证。光计算方面，Lightmatter和Lightelligence等公司则专注于AI推理加速和芯片间互连等具体应用场景，以"加速卡"或"互连芯片"的形态进入现有数据中心。

这种"新兴计算+经典计算"的并行架构，与笔者在2021年提出的观点高度一致：新兴计算在相当长的时期里也将与传统计算通过微服务（API）组成专业计算机。换言之，量子计算和光计算并非要取代硅-冯范式，而是在硅-冯范式难以胜任的特定领域发挥其独特优势。

混合计算的核心意义，在于突破了半导体循环律所依赖的"单一物理机制"假设。传统半导体产业建立在单一的电荷变换物理机制之上——电子在场效应晶体管中的迁移和势垒变化，构成了整个数字世界的物理基础。这一机制的优势在于简单、可靠、可大规模制造，但其局限性也十分明显：电子的运动受制于电阻、热噪声和量子隧穿效应，决定了CMOS微缩的物理极限。

混合计算探索了完全不同的物理机制——光子的无质量传播、量子态的叠加与纠缠、自旋的磁矩取向——这些机制在理论上具有超越电子计算的潜力。摩尔定律成就了两代人的辉煌，也固化了相应的定势思维模式；后摩尔时代呈现百花齐放盛景，新一代人或将创造更大的辉煌。

混合计算成为重要趋势，是半导体循环律终结在物理基础层面的必然延伸。当电子计算的物理极限逐渐显现时，产业必然开始探索基于其他物理机制的计算方式。量子计算和光计算虽然尚处于产业化早期，但已展现出在特定领域超越经典计算的潜力。需要强调的是，混合计算在可预见的未来并非要取代硅-冯范式，而是与其并存，在各自擅长的领域发挥优势。这是后循环时代"多路线并存"逻辑在物理基础层面的展开。

\section{先进封装重构系统边界}

在传统半导体产业中，"芯片"和"系统"之间有清晰的界限：芯片是封装好的单个die，系统是将多颗芯片在PCB上连接起来。随着摩尔定律放缓、单芯片集成难度剧增，这一界限正在被先进封装技术模糊乃至消解。

2019年，笔者在《复归于道——封装改道芯片业》的演讲中提出了一个核心判断："在后摩尔时代，异构/异质集成推动了多芯片封装/多芯片模组的发展，有望在现有芯片产业基础上创造新的产业生态系统与商业模式。"其中特别指出了两个标志性事件：一是2016年ITRS被IRDS取代，宣告摩尔定律驱动的技术路线图终结；二是同年苹果A10处理器采用InFO先进封装技术，标志着封装从"后道工序"走向"系统集成核心"。

先进封装（advanced packaging）正是在这一背景下，从"后道工序"跃升为"系统集成核心"。它的本质是：在封装尺度上实现原本需要在晶圆工艺中完成的集成功能，从而将"系统"的边界从板级推回到封装级。

先进封装在近年来呈现出多条技术路线并行的局面，它们在互连密度、集成方式和应用场景上各有侧重。

TSMC的CoWoS（Chip-on-Wafer-on-Substrate）技术是2.5D集成的代表。通过在硅中介层（silicon interposer）上实现逻辑die与HBM的高密度互连，CoWoS突破了单芯片光罩面积的限制。CoWoS-S Gen5的中介层面积达到约2500mm²（3.3倍光罩拼接），支持8个HBM3堆叠。NVIDIA的H100和B200 GPU均采用这一技术。CoWoS-L则进一步引入局部硅互联桥（LSI）嵌入RDL中介层，兼顾了硅桥的高密度互连与有机基板的尺寸灵活性。

Intel的EMIB（Embedded Multi-die Interconnect Bridge）技术则选择了不同的工程思路：在有机封装基板中嵌入小型硅桥，仅在桥接区域提供高密度互连，从而避免了全硅中介层的高成本。EMIB-T的桥间距达到35μm，能够支持120×180mm的封装尺寸、38+个桥和12+个光罩级die。

「注9」：CoWoS和EMIB代表了2.5D集成的两种工程哲学。CoWoS用大面积硅中介层实现全局高密度互连，EMIB用局部硅桥实现关键区域的高密度互连。前者追求极致性能，后者追求成本效益。

3D集成代表了先进封装中互连密度的最高水平。TSMC的SoIC（System-on-Integrated-Chips）采用无焊料的铜-铜混合键合（hybrid bonding），互连间距达到6μm，密度超过25,000/mm²，带宽密度超过900 Tbps/mm²，能效低于0.05 pJ/bit。AMD的3D V-Cache采用这一技术，将额外64MB L3 SRAM垂直堆叠在CPU核心之上，总L3缓存达到96MB/CCD，而额外延迟仅为4个时钟周期。

Intel的Foveros Direct则将混合键合互连间距进一步压缩到5μm以下，密度超过10,000/mm²。Foveros Direct 3D允许将采用不同制程节点的功能die（计算、I/O、SoC）垂直堆叠成单一"类单片"系统。Meteor Lake处理器首次将CPU计算tile分离到Intel 4制程，而SoC tile和I/O tile使用成熟制程，实现了"异构集成"在三维空间中的展开。

AMD的chiplet架构（自Zen 2起）是芯粒化设计在大规模量产中的成功实践。通过将计算核心（CCD）和I/O（IOD）分离至不同制程节点，AMD实现了"计算核用先进制程、I/O用成熟制程"的优化组合。每颗CCD在TSMC 7nm工艺上制造，而IOD在GF 12nm工艺上制造，通过Infinity Fabric（GMI3）实现缓存一致性互连。这一设计使得AMD能够在单封装内集成最多8个CCD（EPYC Rome，64核），大幅降低了芯片成本。

Apple的UltraFusion技术则展示了硅桥互联的另一种可能。基于TSMC的InFO-LSI技术，UltraFusion通过10,000+个高速信号连接两颗M1 Max die，实现2.5TB/s的片间带宽。macOS将这两颗die视为单颗处理器，在软件层面实现了完全透明的多芯片集成。

「注10」：AMD的chiplet和Apple的UltraFusion代表了chiplet集成的两种不同模式——前者是"异构功能拆分"，后者是"同质算力叠加"。

UCIe（Universal Chiplet Interconnect Express）标准于2022年由AMD、Arm、ASE、Google Cloud、Intel、Meta、Microsoft、Qualcomm、Samsung、TSMC等联合发起，定义了chiplet间互连的开放标准。UCIe 1.0支持2D/2.5D封装，数据率4-32 GT/s，延迟低于2ns。UCIe 2.0于2024年新增了3D堆叠支持和统一安全模型。这一标准的出现，使得不同公司、不同制程、不同功能的chiplet可以在同一个封装内"即插即用"——这本质上是"系统级集成"从专有方案走向开放生态的标志。

HBM（High Bandwidth Memory）的演进则是先进封装在存储领域的典型应用。从HBM1的128 GB/s到HBM3E的1,280 GB/s，再到HBM4预计超过2 TB/s，DRAM堆叠通过TSV和微凸点技术持续突破带宽限制。HBM与逻辑die通过CoWoS/EMIB等2.5D集成技术连接，形成了"计算+存储"在封装级的高效协同。

先进封装对半导体产业最深刻的改变，在于封装的产业定位发生了根本性转变。在传统分工中，封装被视为"后道工序"，是在芯片制造完成后进行的保护性和连接性工作。但在先进封装时代，封装已成为决定系统性能、功耗和尺寸的关键环节。正如笔者在2022年所提出的："封测已经不完全是后道的一个工序了，而是整合整个产业链的部分。从IP、芯片设计开始，就需要封装工艺的配合，将不同功能规划到不同位置，并考虑散热平衡。"

这一转变的产业影响极为深远：封装企业正在从"代工"角色走向"系统集成"角色，芯片设计公司需要从设计之初就考虑封装方案，而EDA工具也需要将封装和芯片设计纳入统一的协同设计框架。

先进封装对半导体循环律的意义，在于它从根本上消解了"统一芯片"范式的物理基础。在统一范式时代，产业的核心假设是"一片晶圆上的所有功能越集成越好"——因为集成意味着更高的性能和更低的成本。但当微缩不再经济时，"单片集成"本身不再是必然的选择。先进封装提供了另一个选项：将多颗较小、较成熟的芯片通过高密度互连组合在一起，实现原本需要在单颗大芯片上才能达到的功能和性能。

这意味着"系统"的边界发生了根本性的移动——从"芯片"推向"封装"，从"die内"推向"die间"。S曲线的视角下，传统单芯片集成已逼近其极限，而基于先进封装的系统级集成正处于快速增长期的起点。

先进封装重构系统边界，是半导体循环律终结在集成方式层面的直接体现。从CoWoS到SoIC，从EMIB到UCIe，从chiplet到HBM，这些技术共同指向一个方向：在后循环时代，"系统"不再是多颗芯片在PCB上的组合，而是在封装层面实现的高密度异构集成体。封装正在从"后道工序"变为"系统集成的核心"，而这一转变的本质，是"从芯片为中心走向系统为中心"这一后循环时代主线在集成技术层面的具象化。

\section{软件定义系统成为新的性能组成方式}

在摩尔定律的黄金时代，硬件性能的持续提升是"免费"的——每一代新工艺都带来晶体管密度翻倍、性能提升、功耗降低。软件开发者只需等待下一代硬件，应用就能自动跑得更快。然而，当登纳德缩放失效、摩尔定律放缓之后，这一"免费午餐"宣告终结。性能提升不能再依赖工艺进步，而越来越依赖于系统级的软硬件协同优化。

由此催生了一个深刻的变化：软件正在从硬件的"使用者"转变为硬件的"定义者"。这就是软件定义系统（software-defined system）的核心思想——通过软件抽象层，将物理硬件资源的组合和调度方式交给软件来控制，从而实现硬件性能的灵活编排和动态适配。

笔者院士在2021年的演讲中对此有过精辟的论述：关于软件定义，可以用两句话概括——"有限的资源，无限的畅想。"他指出，软件定义的核心是API（应用编程接口），它解耦了软、硬件之间的关系：在API之上可以编程，在API之下，基础设施部分没有必要时基本上不动。软件定义有三大特点：推动应用软件向个性化发展，硬件资源向标准化发展，系统功能向智能化发展。

NVIDIA的CUDA平台是软件定义计算最成功的范例。2006年推出的CUDA通过C/C++语言扩展将GPU从固定功能的图形处理器转变为通用的并行计算引擎。开发者通过CUDA编写并行程序，而NVIDIA的驱动和运行时系统负责将这些程序映射到GPU的数千个核心上。这一模型抽象了底层的硬件细节——开发者不需要关心GPU中有多少核心、它们如何连接、缓存如何配置——只需按照CUDA编程模型编写代码，软件自动实现了对硬件性能的组合和调度。

这一模式的成功，使GPU从一个专用图形芯片演变为通用数据中心的算力基石。在2024年，NVIDIA的数据中心业务收入超过300亿美元，证明了"软件定义硬件"的商业价值。

Intel的oneAPI和AMD的ROCm代表了软件定义系统的另一个方向：跨平台硬件抽象层。oneAPI基于SYCL标准，提供对CPU、GPU、FPGA和AI加速器的统一编程接口，开发者编写一次代码即可在不同硬件后端上运行。ROCm的HIP（Heterogeneous Interface for Portability）则允许一套C++代码同时在AMD和NVIDIA GPU上运行。

这些抽象层的本质，是将硬件差异通过软件"抹平"——在应用开发者看来，不同厂商、不同架构的硬件只是同一抽象层的不同实现。这正是软件定义系统的核心特征：用软件去定义和组合硬件能力。

AMD Xilinx的Vitis统一软件平台代表了软件定义系统的另一个维度：让原本需要硬件专业知识才能使用的FPGA，通过软件编程即可完成硬件加速。Vitis支持C、C++、OpenCL和Python来编写加速器，通过高层次综合（HLS）技术将高级语言代码自动编译为FPGA比特流。AWS EC2 F1 FPGA实例则进一步将这种能力以云服务的形式交付——用户只需上传代码，就能在云端获得可编程的硬件加速资源，而无需自行管理FPGA硬件。

``软件定义一切的三个层次''——笔者在2021年提出，可以从器件层、架构层和系统层三个维度来看待这一趋势。在器件层，软件定义硬件（SDH）使芯片可以像CPU一样编程，同时具备接近ASIC的性能；在架构层，软件定义架构（SDA）将硬件和软件分隔开来，重塑了信息基础架构；在系统层，API解耦了软硬件之间的关系，使硬件资源向标准化方向发展，应用软件向个性化方向发展。

RISC-V的开放指令集架构为软件定义系统提供了全新的维度——软件不仅定义了硬件的功能，还定义了硬件的指令集本身。RISC-V在设计之初就预留了custom-0到custom-3的编码空间，芯片设计者可以根据应用需求添加自定义指令。这意味着，对于一个特定的计算任务，软件（编译器）和硬件（指令扩展）可以协同优化——软件定义了需要什么样的硬件支持，硬件实现这些支持后反过来使软件运行得更快。

DARPA的软件定义硬件（Software-Defined Hardware，SDH）项目试图将软件定义推向极致：构建支持运行时重构的硬件和软件，使数据密集型算法在不牺牲可编程性的前提下实现接近ASIC的性能。其目标是达到ASIC效率的5倍以内，比CPU实现好500-1000倍，同时保持与Python等同的可编程性。如果这一目标得以实现，它将彻底消除"通用"和"专用"之间的鸿沟。

「注11」：DARPA的SDH项目于2017年作为"电子复兴计划"（ERI）的一部分启动，Intel、NVIDIA、Qualcomm、Georgia Tech和Stanford等均为承研单位。这一项目的背景与笔者所关注的"后摩尔时代架构创新"高度一致。

Tesla的Dojo超算是软件定义系统最极端的实践：它完全从FSD（全自动驾驶）神经网络的训练需求出发，反向定义了芯片架构和系统拓扑。Dojo的D1芯片包含354个完全可编程的处理节点，采用自定义TTPoE（Tesla Transport Protocol over Ethernet）协议实现软件定义网络，使用纯2D mesh拓扑而非传统的GPU层次化互联结构。

Dojo没有DRAM，每个计算节点只有小容量SRAM——因为FSD网络的权值可以通过数据流的方式实时加载。这一设计完全突破了传统处理器"大缓存+大内存"的架构范式，转而追求"持续的数据流处理"。Dojo的每一个设计决策，都源于特定软件需求的直接投射——这是软件定义系统的终极形态：软件不仅调度硬件，还塑造了硬件的每一个细节。

软件定义系统对半导体产业的影响，正在催生一种新的商业模式：硬件即服务（Hardware-as-a-Service）。在传统模式下，客户购买硬件，自行编写软件来使用它。在软件定义模式下，硬件能力通过API以服务的形式交付，客户不再关心底层硬件是什么、在哪里、如何配置——他们只需要调用API，软件会自动完成所有硬件资源的编排。

这种模式的典型代表是云计算中的FPGA实例（AWS F1）和GPU实例（NVIDIA GPU Cloud）。在这些平台上，硬件的选择、配置、调度完全由软件层完成，用户获得的是"算力"而非"芯片"。这改变了半导体产业的价值链：从"卖芯片"转向"卖算力"。

笔者敏锐地观察到了这一趋势，他在2021年的演讲中指出："现在我们通过API连接异构计算，未来新兴计算在相当长的时期里也将与传统计算通过微服务（API）组成专业计算机。"

软件定义系统的兴起，是半导体循环律"统一"逻辑从硬件层向软件层延伸的体现。在传统统一芯片范式中，"统一"的实现方式是硬件层面的统一——统一的CMOS工艺、统一的冯·诺依曼架构、统一的x86指令集。在后循环时代，"统一"的方式正在从硬件转向软件——API成为新的"统一抽象层"，硬件资源被软件层封装和标准化，通过编程接口向上层应用提供服务。

这一转变的深层意义在于：当硬件层面的"统一"因微缩极限和架构发散而瓦解后，软件层面的"统一"成为维持系统级协同的新途径。从笔者的"由简入繁、化繁为简"史观来看，这是产业在经历"硬件分化"（由简入繁）之后，通过"软件统一"（化繁为简）实现螺旋式上升的典型表现。

软件定义系统成为新的性能组成方式，是半导体循环律终结后统一逻辑在软件层面的延伸和升华。从CUDA到oneAPI，从Vitis到RISC-V，从DARPA SDH到Tesla Dojo，这一系列实践表明：当硬件层面的"免费午餐"结束后，性能提升越来越多地来源于软件对硬件能力的灵活编排和动态适配。软件正在从"使用者"变为"定义者"——这是后循环时代"系统协同"逻辑在软件层面的集中体现。正如笔者所总结的："有限的资源，无限的畅想。"

% 第四章结束

\end{document}
